\section{Characteristic lengths of the region births}
\label{sec:area_law_repair_lengths}

% Source: 10-geometry.tex, geometry:primary-scale, lines 709--721;
% belt and replacement scanner scales, lines 732--740;
% old and fresh identifiers, lines 398--409; birth order, lines 565--568.
% Revision: adc7f1241b42e322a6451854ab7e4b4c146bf78a.
\begin{definition}[Characteristic birth lengths]
    \label{def:area_law_repair_characteristic_length}
    \lean{TNLean.PEPS.AreaLaw.Geometry.initialRegionCharacteristicLength,
        TNLean.PEPS.AreaLaw.Geometry.repairStageCharacteristicLength}
    \leanok
    \uses{def:area_law_repair_stage}
    For an initial identifier family, define its characteristic length
    $\tau_0:I_0\to\mathbb R$ by
    \begin{align}
        \tau_0(i) & =
        \begin{cases}
            0,   & i\text{ is the dummy},                         \\
            r_k, & i\text{ is a primary in layer }k,              \\
            t_k, & i\text{ is an initial belt region in layer }k.
        \end{cases}
        \label{eq:area_law_initial_characteristic_length}
    \end{align}
    Here $r_k=2^k$ is the layer-cell side and $t_k$ is the fine-cell
    side in layer $k$. The primary pitch side $s_k$ can be larger than
    $r_k$; the characteristic length uses $r_k$.

    Let $\mathcal D_j$ be the recursive family of
    Definition~\ref{def:area_law_repair_stage}, with identifier set
    $I_j$. Start with the assignment above for its jointly chosen
    initial family. Identify each old $i\in I_j$ with its retained
    copy in $I_{j+1}$. A fresh identifier $q=(j+1,p,H)$ records its
    parent mark $p$ and run $H$; write $S_p^{(j)}$ for that parent's
    scale at stage $j$. For each old identifier and each fresh
    identifier, respectively, define
    \begin{align}
        \tau_{j+1}(i)       & =\tau_j(i),    \\
        \tau_{j+1}(j+1,p,H) & =d(S_p^{(j)}).
        \label{eq:area_law_repair_characteristic_length}
    \end{align}
\end{definition}

A replacement with characteristic length $u$ has removal half-side
$u$ and scanner primitive scale $2u$, as in the scanner-scale
assignment of \cite[Section~11]{OpenAI2026AreaLaw}.

The dummy is omitted from the retained family, so its assigned
length will not be used. Birth order places primaries
first, initial belt regions second, and replacements by generation,
with arbitrary ties within each class.

% Source: 10-geometry.tex, retention and creation, lines 398--409;
% characteristic lengths and scanner scales, lines 709--740;
% geometry:total-repairs, lines 685--692; nonprimary errors, lines 796--798.
% Revision: adc7f1241b42e322a6451854ab7e4b4c146bf78a.

\begin{theorem}[Positive characteristic lengths]
    \label{thm:area_law_repair_characteristic_length_positive}
    \lean{TNLean.PEPS.AreaLaw.Geometry.repairStageCharacteristicLength_pos_of_ne_dummy}
    \leanok
    \uses{def:area_law_repair_characteristic_length,def:area_law_repair_stage}
    Fix an origin $o\in\mathbb R^2$, a finite domain
    $\Lambda\subset\mathbb Z^2$, a cut $A\subseteq\Lambda$, integers
    $C_0\ge2$ and $k_0\ge50{,}000{,}000$, and any $S_*\in\mathbb R$.
    In the actual repair family, every non-dummy identifier at every
    generation has positive characteristic length:
    \begin{align}
        \tau_N(i) & >0
        \qquad(N\in\mathbb N,\ i\in I_N,\ i\ne *_N).\notag
    \end{align}
\end{theorem}

\begin{proof}
    \leanok
    \uses{def:area_law_repair_characteristic_length,
        def:area_law_repair_square,def:area_law_repair_stage}
    Induct on the generation. Initially, primary and belt lengths
    are positive powers of two. An old non-dummy identifier retains
    its length and remains distinct from the old dummy. A fresh
    identifier has length $d(S)$, which is a power of two with an
    integer exponent and hence positive.
\end{proof}

% Source: 10-geometry.tex, lines 212--218, 398--409, 497--504,
% 558--571, 709--721 and 732--740.
% Revision: adc7f1241b42e322a6451854ab7e4b4c146bf78a.

\begin{lemma}[Preservation under retention]
    \label{lem:area_law_repair_retention_identities}
    \lean{TNLean.PEPS.AreaLaw.Geometry.repairStageOldInclusion_injective,
        TNLean.PEPS.AreaLaw.Geometry.repairStageOldInclusion_dummy,
        TNLean.PEPS.AreaLaw.Geometry.repairStage_color_oldInclusion,
        TNLean.PEPS.AreaLaw.Geometry.repairStageCharacteristicLength_oldInclusion}
    \leanok
    \uses{def:area_law_repair_inclusions,def:area_law_repair_characteristic_length,
        thm:area_law_retained_birth_lattice_point}
    Fix an origin $o\in\mathbb R^2$, a finite domain
    $\Lambda\subset\mathbb Z^2$, a cut $A\subseteq\Lambda$, integers
    $C_0\ge2$ and $k_0\ge50{,}000{,}000$, and any $S_*\in\mathbb R$.
    Consider the actual repair family $\mathcal D_j$, with its jointly
    chosen initial residues. For every $j,n\in\mathbb N$, the retention
    map $\iota_{j,j+n}:I_j\to I_{j+n}$ is injective and satisfies
    $\iota_{j,j+n}(*_j)=*_{j+n}$. For every $i\in I_j$,
    \begin{align}
        B^{(j+n)}_{\iota_{j,j+n}(i)}   & =B^{(j)}_i,\notag   \\
        \xi^{(j+n)}_{\iota_{j,j+n}(i)} & =\xi^{(j)}_i,\notag \\
        \tau_{j+n}(\iota_{j,j+n}(i))   & =\tau_j(i).\notag
    \end{align}
    Thus retention preserves the remembered birth, its binary label,
    and its characteristic length, as in the repair construction of
    \cite[Section~11]{OpenAI2026AreaLaw}.
\end{lemma}

\begin{proof}
    \leanok
    \uses{def:area_law_repair_inclusions,def:area_law_repair_stage,
        def:area_law_repair_successor_regions,
        def:area_law_repair_characteristic_length,
        thm:area_law_retained_birth_lattice_point}
    Induct on $n$. For $n=0$ the retention map is the identity.
    Each subsequent step includes the old identifier family into a
    disjoint union, retains its dummy, and leaves all three recorded
    quantities unchanged. Composition preserves injectivity and the
    displayed identities.
\end{proof}

\begin{lemma}[Rank-zero ancestors at an earlier generation]
    \label{lem:area_law_repair_rank_zero_ancestry}
    \lean{TNLean.PEPS.AreaLaw.Geometry.exists_repairStageOldInclusion_of_repairStageBirthRank_eq_zero}
    \leanok
    \uses{def:area_law_repair_inclusions,def:area_law_repair_birth_rank}
    For the same actual repair family, let $j,n\in\mathbb N$ and
    $x\in I_{j+n}$ satisfy $\operatorname{rk}_{j+n}(x)=0$.
    There is $y\in I_j$ such that
    $\operatorname{rk}_j(y)=0$ and $\iota_{j,j+n}(y)=x$.
    The assertion includes the dummy identifier.
\end{lemma}

\begin{proof}
    \leanok
    \uses{def:area_law_repair_inclusions,def:area_law_repair_birth_rank}
    Induct on $n$. The case $n=0$ follows by taking $y=x$.
    At generation $j+n+1$, a fresh identifier has rank $j+n+2>0$.
    A rank-zero identifier is therefore old, and the induction
    hypothesis gives its ancestor at generation $j$.
\end{proof}

\begin{lemma}[The unique original primary of a rank-zero identifier]
    \label{lem:area_law_repair_rank_zero_primary_ancestor}
    \lean{TNLean.PEPS.AreaLaw.Geometry.existsUnique_initialPrimary_of_repairStageBirthRank_eq_zero}
    \leanok
    \uses{def:area_law_repair_inclusions,def:area_law_repair_birth_rank,
        def:area_law_initial_region_indices,def:area_law_primary_pitch_indices}
    Use the same jointly chosen residues as the actual initial stage,
    and let $G_k$ be its primary pitch index set at layer $k$.
    For every $n\in\mathbb N$ and every $x\in I_n$ with
    $\operatorname{rk}_n(x)=0$ and $x\ne *_n$, there is exactly one
    pair $(k,J)$ with $k\ge k_0$ and $J\in G_k$ such that
    $\iota_{0,n}((k,J))=x$.
    Here $(k,J)$ is the original primary identifier in $I_0$.
    No finiteness of the family over all layers is assumed.
\end{lemma}

\begin{proof}
    \leanok
    \uses{lem:area_law_repair_rank_zero_ancestry,
        lem:area_law_repair_retention_identities,
        def:area_law_initial_region_indices,def:area_law_repair_birth_rank}
    Take a rank-zero ancestor in $I_0$. It cannot be the dummy,
    since retention maps the dummy to $*_n$, and it cannot be an
    initial belt identifier, whose rank is one. It is therefore
    an original primary $(k,J)$. Injectivity of retention and of
    the primary summand gives uniqueness of this pair.
\end{proof}
